\SourcePage{699}{702}
\section{Scattering with spin--orbit interaction}

So far we have considered only collisions of particles whose interaction is
independent of their spins. Under these conditions the spins either have no
effect on scattering at all or act indirectly through exchange effects
(\S~137).

We now generalize the scattering theory developed in \S~123 to the case in
which the interaction depends essentially on the particle spins, as in
collisions of nuclear particles.

Consider in detail the simplest case, where one colliding particle, taken for
definiteness to be the incident-beam particle, has spin $1/2$, and the other,
the target particle, has spin zero.

For a specified half-integral total angular momentum $j$, the orbital angular
momentum can have only the two values $l=j\pm1/2$, corresponding to states of
different parity. Conservation of $j$ and parity therefore also implies
conservation of the magnitude of the orbital angular momentum.

The operator $\hat f$ (see \S~125) now acts on both the orbital and spin
variables of the system's wave function. It must commute with the operator of
the conserved quantity $\hat{\mathbf l}^{,2}$. Its most general form is
\begin{equation}
 \hat f=\hat a+\hat b\,\hat{\mathbf l}\cdot\hat{\mathbf s},
 \tag{140.1}
\end{equation}
where $\hat a$, $\hat b$ are orbital operators depending only on
$\hat{\mathbf l}^{,2}$.

The $S$-matrix, and hence the matrix of $\hat f$, is diagonal with respect to
the wave functions of states having specified values of the conserved
quantities $l$ and $j$ (and of the projection $m$ of the total angular
momentum); its diagonal elements are expressed in terms of the wave-function
phases $\delta$ by (123.15). At fixed $l$, for total angular momentum
$j=l+1/2$ and $j=l-1/2$, the eigenvalues of
$\hat{\mathbf l}\cdot\hat{\mathbf s}$ are respectively $l/2$ and
$-(l+1)/2$ (see (118.5)). Thus the diagonal matrix elements $a_l$ and $b_l$
of $\hat a$ and $\hat b$ obey
\begin{equation}
 a_l+\frac l2b_l=\frac1{2ik}(e^{2i\delta_l^+}-1),\qquad
 a_l-\frac{l+1}{2}b_l=\frac1{2ik}(e^{2i\delta_l^-}-1),
 \tag{140.2}
\end{equation}
where $\delta_l^+$ and $\delta_l^-$ correspond to the states $j=l+1/2$ and
$j=l-1/2$.

Our interest, however, is not in the diagonal elements of $\hat f$ themselves
for states of specified $l$

\SourcePage{700}{703}
and $j$, but in the scattering amplitude as a function of the directions of
the incident and scattered waves. This amplitude is still an operator, but
only in the spin variables; it is not diagonal in the spin projection
$\sigma$. In the remainder of this section, $\hat f$ denotes this operator.

To find it, apply (140.1) to the function (125.17) corresponding to a plane
wave incident along the $z$ axis. Thus
\begin{equation}
 \hat f=\sum_{l=0}^{\infty}(2l+1)(a_l+b_l\hat{\mathbf l}\cdot\hat{\mathbf s})P_l(\cos\theta).
 \tag{140.3}
\end{equation}
It remains to evaluate the action of
$\hat{\mathbf l}\cdot\hat{\mathbf s}$ on $P_l(\cos\theta)$. This may be done
by writing
\[
 \hat{\mathbf l}\cdot\hat{\mathbf s}
 =\frac12(\hat l_+\hat s_-+\hat l_-\hat s_+)+\hat l_z\hat s_z
\]
(see (29.11)) and using (27.12) for the matrix elements of $l_\pm$; it is even
simpler to use the operator expressions (26.14), (26.15) directly. A short
calculation gives
\[
 \hat{\mathbf l}\cdot\hat{\mathbf s}P_l(\cos\theta)
 =i\mathbf v\cdot\hat{\mathbf s}P_l^1(\cos\theta),
\]
where $P_l^1$ is an associated Legendre polynomial and $\mathbf v$ is a unit
vector in the direction $\mathbf n\times\mathbf n'$, perpendicular to the
scattering plane. Here $\mathbf n$ is the incident direction along the $z$
axis and $\mathbf n'$ is the scattering direction specified by the spherical
angles $\theta$, $\varphi$.

Solving (140.2) for $a_l$, $b_l$ and substituting in (140.3), we obtain
\begin{equation}
 \hat f=A+2B\mathbf v\cdot\hat{\mathbf s},
 \tag{140.4}
\end{equation}
\begin{align}
 A&=\frac1{2ik}\sum_{l=0}^{\infty}[(l+1)(e^{2i\delta_l^+}-1)+l(e^{2i\delta_l^-}-1)]P_l(\cos\theta),\notag\\
 B&=\frac1{2k}\sum_{l=1}^{\infty}(e^{2i\delta_l^+}-e^{2i\delta_l^-})P_l^1(\cos\theta).
 \tag{140.5}
\end{align}
The matrix elements of this operator are scattering amplitudes for specified
values of the spin projection in the initial ($\sigma$) and final
($\sigma'$) states. Consider a cross-section summed over every possible
$\sigma'$ and averaged with the probabilities of the different $\sigma$ in
the initial state, that is, in the

\SourcePage{701}{704}
incident beam. It is calculated as
\begin{equation}
 d\sigma=\overline{(\hat f^+\hat f)_{\sigma\sigma}}d\Omega;
 \tag{140.6}
\end{equation}
taking diagonal elements of $\hat f^+\hat f$ performs the sum over final
states, while the bar denotes averaging over the initial
state\footnote{If the squared modulus $|f_{0n}|^2$ of a matrix element of an
operator for a transition $0\to n$ is summed over the final states $n$, then
\[
 \sum_n|f_{0n}|^2=\sum_nf_{0n}(f_{0n})^*=\sum_nf_{0n}(f^+)_{n0}=(ff^+)_{00}.
\]
To avoid misunderstanding, we emphasize that the adjoint sign $+$ in (140.6)
and below refers to $\hat f$ as a spin operator and, in particular, does not
imply interchange of $\mathbf n$ and $\mathbf n'$.}. If all spin directions
are equally probable in the initial state, this average reduces to taking the
trace of the matrix divided by the number of possible spin projections:
\begin{equation}
 d\sigma=\frac12\Tr(\hat f^+\hat f)d\Omega.
 \tag{140.7}
\end{equation}

On substituting (140.4) in (140.6), the mean square
$(\mathbf v\cdot\mathbf s)^2$ is
$v^2s^2/3=s(s+1)/3=1/4$. The result is
\begin{equation}
 \frac{d\sigma}{d\Omega}=|A|^2+|B|^2
 +2\operatorname{Re}(AB^*)\mathbf v\cdot\mathbf P,
 \tag{140.8}
\end{equation}
where $\mathbf P=2\bar{\mathbf s}$ is the initial beam polarization, defined
as the ratio of the mean initial spin to its greatest possible value
($1/2$). Recall that for spin $1/2$, the vector $\bar{\mathbf s}$ completely
characterizes the spin state (\S~59).

Notice that polarization of the incident beam produces an azimuthal asymmetry
of the scattering. Through the factor $\mathbf v\cdot\mathbf P$ in the last
term, (140.8) depends not only on the polar angle $\theta$ but also on the
azimuth $\varphi$ of $\mathbf n'$ relative to $\mathbf n$ (provided the
polarization is not perpendicular to $\mathbf v$, so that
$\mathbf v\cdot\mathbf P\ne0$).

The polarization of the scattered particles can be calculated from
\begin{equation}
 \mathbf P'=2\frac{\overline{(\hat f^+\mathbf s\hat f)_{\sigma\sigma}}}
 {\overline{(\hat f^+\hat f)_{\sigma\sigma}}}.
 \tag{140.9}
\end{equation}
For example, if the initial state is unpolarized ($\mathbf P=0$), a short
calculation gives
\begin{equation}
 \mathbf P'=\frac{2\operatorname{Re}(AB^*)\mathbf v}{|A|^2+|B|^2}.
 \tag{140.10}
\end{equation}

\SourcePage{702}{705}
Thus scattering generally produces polarization perpendicular to the
scattering plane. This effect is absent in the Born approximation, however:
if all phases $\delta$ are small, then to first order in them $A$ is real and
$B$ is purely imaginary, so $\operatorname{Re}(AB^*)=0$.

That $\mathbf P'$ in (140.10) is directed along $\mathbf v$ is evident in
advance: $\mathbf P'$ is an axial vector, and $\mathbf v$ is the only axial
vector that can be formed from the available polar vectors $\mathbf n$ and
$\mathbf n'$. The same property therefore also holds for the polarization
produced when an unpolarized beam of spin-$1/2$ particles is scattered by an
unpolarized target of nuclei with any spin, not only zero\footnote{Here we
mean a target of nuclei whose spin directions are completely random. Recall
that for $s>1/2$ the mean spin vector does not fully determine the spin state;
its vanishing does not by itself mean a complete absence of spin ordering.}.

The formulation of the reciprocity theorem for scattering in the presence of
spins must take account of the fact that time reversal changes the signs not
only of momenta but also of angular momenta. Symmetry under time reversal must
therefore be expressed as equality of the amplitudes of processes differing
not only by interchange of initial and final states and reversal of the
directions of motion, but also by reversal of the spin projections in both
states. The signs of these amplitudes may nevertheless differ, since time
reversal introduces into the spin wave function the factor
$(-1)^{s-\sigma}$ according to (60.3). The reciprocity theorem must therefore
be formulated as follows\footnote{The derivation is analogous to that of
(125.12). Spin factors must be introduced in the amplitudes of the incoming
and outgoing waves in the wave function, and (125.10) is replaced by
$\hat K^{-1}\hat S\hat K=\hat S$, where $\hat K$ not only performs inversion
but also transforms the spin state according to (60.3).}:
\begin{align}
 &f(\sigma_1,\sigma_2,\mathbf n;\sigma_1',\sigma_2',\mathbf n')\notag\\
 &\quad=(-1)^{\sum(s-\sigma)}
 f(-\sigma_1',-\sigma_2',-\mathbf n';-\sigma_1,-\sigma_2,-\mathbf n).
 \tag{140.11}
\end{align}
Here $f(\sigma_1,\sigma_2,\mathbf n;\sigma_1',\sigma_2',\mathbf n')$ is the
scattering amplitude for the spin projections of the colliding particles to
change from $\sigma_1$, $\sigma_2$ to $\sigma_1'$, $\sigma_2'$. The sum in
the exponent is over both particles before and after the scattering.

\SourcePage{703}{706}
In the Born approximation scattering has an additional symmetry: the
probabilities are the same for processes differing by interchange of the
initial and final states without reversal of the particle momenta and spin
projections, as under time reversal (see \S~126). Combining this property with
the reciprocity theorem, we find that scattering is symmetric under reversal
of all momenta and spin projections without their interchange. It follows
that, in the Born approximation, polarization cannot be produced when any
unpolarized beam is scattered by an unpolarized target. Indeed, under this
transformation the polarization vector $\mathbf P$ changes sign, whereas the
only vector $\mathbf k\times\mathbf k'$ along which $\mathbf P$ can be
directed is unchanged. Thus the property noted above for spin-$1/2$ particles
scattered by spin-zero particles is in fact general.

For arbitrary spins of the colliding particles, the general formulas for the
angular distributions are very cumbersome, and we shall not derive them
here. We shall only count the parameters that determine these distributions.

The case considered above, of particles with spins $1/2$ and zero, is
characterized in particular by there being only one state of the two-particle
system for specified values of $J$ and parity, apart from the immaterial
orientation of the total angular momentum in space. Each such state
contributes one real parameter, a phase $\delta$, to the scattering amplitude.
For other spins, there are generally several different states with the same
total angular momentum $J$ and parity; they differ in the total particle spin
$S$ and the orbital angular momentum $l$ of their relative motion. Let the
number of these states be $n$. It is readily seen that every such group
contributes $n(n+1)/2$ independent real parameters to the scattering
amplitude.

Indeed, for these states the $S$-matrix is a unitary symmetric matrix, by the
reciprocity theorem, with $n\cdot n$ complex elements. The number of
independent quantities is conveniently counted by writing
$\hat S=\exp(i\hat R)$. Unitarity is then automatic when $\hat R$ is an
arbitrary Hermitian operator (see (12.13)). If $S$ is symmetric, so is $R$;
being Hermitian, it is then real. A real symmetric matrix

\SourcePage{704}{707}
has $n(n+1)/2$ independent components.

For example, for two spin-$1/2$ particles $n=2$. Indeed, at fixed $J$ there
are four states: two with $l=J$ and total spin $S=0$ or 1, and two with
$l=J\pm1$, $S=1$. Clearly two are even ($l$ even) and two are odd ($l$ odd).

The general scattering amplitude for spin-$1/2$ particles, as an operator in
the spin variables of both particles, is readily written from the required
invariance conditions: it must be a scalar invariant under time reversal. We
have available the two axial spin vectors $\mathbf s_1$, $\mathbf s_2$ and
the two ordinary polar vectors $\mathbf n$, $\mathbf n'$. Each of the
operators $\mathbf s_1$, $\mathbf s_2$ must enter the amplitude linearly,
since every function of a spin-$1/2$ operator reduces to a linear function.
The most general operator satisfying these conditions can be written
\begin{align}
 \hat f={}&A+B(\mathbf s_1\cdot\boldsymbol\lambda)(\mathbf s_2\cdot\boldsymbol\lambda)
 +C(\mathbf s_1\cdot\boldsymbol\mu)(\mathbf s_2\cdot\boldsymbol\mu)
 +D(\mathbf s_1\cdot\boldsymbol\nu)(\mathbf s_2\cdot\boldsymbol\nu)\notag\\
 &+E(\mathbf s_1+\mathbf s_2)\cdot\boldsymbol\nu
 +F(\mathbf s_1-\mathbf s_2)\cdot\boldsymbol\nu.
 \tag{140.12}
\end{align}
The coefficients $A$, $B$, \ldots are scalars depending only on
$\mathbf n\cdot\mathbf n'$, that is, on the scattering angle $\theta$ and
the energy. The vectors $\boldsymbol\lambda$, $\boldsymbol\mu$,
$\boldsymbol\nu$ are three mutually perpendicular unit vectors directed
respectively along $\mathbf n+\mathbf n'$, $\mathbf n-\mathbf n'$, and
$\mathbf n\times\mathbf n'$. Time reversal corresponds to
\[
 \mathbf s_1\to-\mathbf s_1,\quad\mathbf s_2\to-\mathbf s_2,\quad
 \mathbf n\to-\mathbf n',\quad\mathbf n'\to-\mathbf n.
\]
Then $\boldsymbol\lambda\to-\boldsymbol\lambda$,
$\boldsymbol\mu\to\boldsymbol\mu$, and
$\boldsymbol\nu\to-\boldsymbol\nu$, making the invariance of (140.12)
evident.

For mutual scattering of nucleons, protons and neutrons, the last term in
(140.12) is absent. This follows already from the fact that the nuclear forces
between nucleons conserve the magnitude of the total spin $S$, whereas
$\mathbf s_1-\mathbf s_2$ does not commute with $S^2$. The remaining terms in
(140.12) can be expressed through the total-spin operator $\mathbf S$ by
(117.4), and therefore commute with $S^2$. For scattering of identical
nucleons, $pp$ or $nn$, the coefficients $A$, $B$, \ldots as functions of the
scattering angle also satisfy symmetry relations following from the identity
of the particles (see Problem 2).

\SourcePage{705}{708}
\subsection*{Problems}

\ProblemHeading 1. For scattering of spin-$1/2$ particles by spin-zero
particles, determine the polarization after scattering when it was also
nonzero before scattering.

\SolutionHeading The calculation from (140.9) is conveniently performed in
components, choosing the $z$ axis along $\mathbf v$. The result is
\[
 \mathbf P'=\frac{(|A|^2-|B|^2)\mathbf P+2|B|^2\mathbf v(\mathbf v\cdot\mathbf P)
 +2\operatorname{Im}(AB^*)\mathbf v\times\mathbf P
 +2\mathbf v\operatorname{Re}(AB^*)}
 {|A|^2+|B|^2+2\operatorname{Re}(AB^*)\mathbf v\cdot\mathbf P}.
\]

\ProblemHeading 2. Find the symmetry conditions satisfied, as functions of
the angle $\theta$, by the coefficients in the scattering amplitude of two
identical nucleons (R.~Oehme, 1955).

\SolutionHeading Regroup the terms in (140.12) so that each is nonzero only
for singlet ($S=0$) or triplet ($S=1$) states of the nucleon system:
\begin{align}
 f={}&a(\mathbf s_1\cdot\mathbf s_2-1/4)+b(\mathbf s_1\cdot\mathbf s_2+3/4)
 +c[1/4+(\mathbf s_1\cdot\boldsymbol\nu)(\mathbf s_2\cdot\boldsymbol\nu)]\notag\\
 &+d[(\mathbf s_1\cdot\mathbf n)(\mathbf s_2\cdot\mathbf n')
 +(\mathbf s_1\cdot\mathbf n')(\mathbf s_2\cdot\mathbf n)]
 +e(\mathbf s_1+\mathbf s_2)\cdot\boldsymbol\nu.
 \tag{1}
\end{align}
Using (117.4), it is readily verified that the first term is nonzero only for
$S=0$, and the others only for $S=1$. By particle identity, the scattering
amplitude must be symmetric under interchange of particle coordinates for
$S=0$ and antisymmetric for $S=1$. This transformation means
$\theta\to\pi-\theta$, or equivalently reversal of one of $\mathbf n$,
$\mathbf n'$ (compare \S~137). The conditions are
\begin{gather}
 a(\pi-\theta)=a(\theta),\quad b(\pi-\theta)=-b(\theta),\quad
 c(\pi-\theta)=-c(\theta),\notag\\
 d(\pi-\theta)=d(\theta),\qquad e(\pi-\theta)=e(\theta).
 \tag{2}
\end{gather}
By isotopic invariance, the scattering amplitude is the same for $nn$ and
$pp$ scattering and for $np$ scattering in the isotopic state $T=1$. The
$np$ system can also be in a state with $T=0$; its scattering amplitude is
then characterized by different coefficients $a$, $b$, \ldots in (1), which
do not have the symmetry properties (2).
